% Quarantine-only candidate statement.
% Date: 29 August 2026.
% This file is not a live-book input and makes no novelty or priority claim.

\begin{lemma}
\label{lem:p213-normal-curve-nonrigid-localization}
Let $K$ be a complete non-Archimedean field, possibly trivially valued,
let $Y$ be a good separated normal $K$-analytic space that is purely
one-dimensional, let $y\in Y$ be nonrigid, and let $W$ be an open
neighborhood of $y$.  There are a compact analytic neighborhood
$V\subset W$ of $y$ and a moment
\[
             F:V\longrightarrow\mathbf G_{m,K}^{N}
\]
such that
\begin{equation}
\label{eq:p213-normal-curve-nonrigid-boundary}
 F_{\mathrm{trop}}(y)\notin
 F_{\mathrm{trop}}\bigl(\partial(V/W)\bigr).
\end{equation}
Consequently, if
\[
 \Omega:=\mathbf R^N\setminus
 F_{\mathrm{trop}}\bigl(\partial(V/W)\bigr),
 \qquad
 U:=F_{\mathrm{trop}}^{-1}(\Omega),
\]
then $U$ is an open neighborhood of $y$ in $W$ and
$F_{\mathrm{trop}}|_U:U\to\Omega$ is proper.
\end{lemma}

\begin{remark}
\label{rem:p213-normal-curve-nonrigid-localization-scope}
Neither strictness nor nontriviality of the ground-field valuation is used.
Normality and nonrigidity enter only in proving that a holomorphic separator
of $y$ from a boundary point cannot vanish at $y$.  The corresponding
assertion at a rigid point is false over a trivially valued field, as shown
by \cite[Remark~12.3.7]{CLD25}.
\end{remark}

